%% ============================================================
%%  Principia Orthogona, Volume Two
%%  Section 2: Contact Realization of the Fold Operator
%% ============================================================

\section{Contact Realization of the Fold Operator}
\label{sec:contact-realization}

This section constructs an explicit contact-geometric realization
of the fold operator $F$ introduced abstractly in Volume
One~\cite{PrincipiaI}, and shows how it arises naturally within
the dm3 framework of Generative Contact
Mechanics~\cite{Paper1,Paper2}. The goal is not to re-derive
the singularity-theoretic properties of folding, but to demonstrate
how the geometric fold appears as a contact-Hamiltonian
discontinuity whose regularized dynamics are governed by the dm3
equations.

Throughout, we assume the existence of the operator sequence
\[
  C \;\to\; K \;\to\; F \;\to\; U
\]
acting on trajectories in a Riemannian manifold $(X,g)$, as
established in Volume One~\cite{PrincipiaI}. In particular, we
take as given: the curvature threshold $\kappa^*$ defined by the
focal radius and corrected by the Rauch comparison theorem; the
fold operator $F$ producing rank-1 loss of injectivity at
$\kappa^*$; and the post-fold stabilization map $U$ defined by
gradient descent on a Morse stability functional. The purpose of
the present work is not to re-derive these structures, but to
construct their contact-geometric realization and to show that
they are dynamically stable within the dm3 framework.

%% ------------------------------------------------------------
\subsection{From Hamiltonian Phase Space to Contact Extension}
\label{subsec:contactification}

In Volume One, the fold is realized in the Hamiltonian phase
space $T^*X$ as a symplectic discontinuity produced by a
distributional term in the action functional. At a fold point
$s_0$, the momentum undergoes an impulsive jump
\[
  p(s_0^+) - p(s_0^-) = \mu\,n(s_0),
\]
where $n(s_0)$ is the unit normal to the compressed trajectory
at the critical curvature point. The fold map
\[
  \mathcal{F} : (\gamma,p) \mapsto (\gamma, p + \mu n)
\]
preserves the canonical symplectic form $\omega = d\gamma\wedge dp$
(Volume One, Theorem on Symplectic Preservation).

However, the Hamiltonian framework alone cannot describe the
post-fold attracting dynamics selected by the unfolding operator
$U$: symplectic flows on compact phase spaces are volume-preserving
and therefore forbid attractors. To capture both the impulsive
fold and the subsequent stabilization in a unified geometric
structure, we pass to the \emph{contact extension}
\[
  M \;=\; X \times \mathbb{R},
\]
with coordinate $z\in\mathbb{R}$ (the action variable) and
contact form
\[
  \alpha \;=\; dz - \lambda,
\]
where $\lambda\in\Omega^1(X)$ satisfies $d\lambda=\omega$. This
is the contactification construction of~\cite{Paper1}, \S6.

%% ------------------------------------------------------------
\subsection{The Fold as a Contact Discontinuity}
\label{subsec:contact-fold}

Within the contact manifold $(M,\alpha)$, the Hamiltonian fold
of Volume One appears as a singular limit of a smooth contact
Hamiltonian vector field. The precise statement is as follows.

The distributional generator of the fold,
\[
  S(\gamma) \;=\; \mu\,\Theta\!\bigl(\kappa(\gamma)-\kappa^*\bigr),
\]
where $\Theta$ is the Heaviside function, induces the momentum
jump $p^+=p^-+\partial S/\partial\gamma$ in the variational
sense (Volume One, \S12). In the contact extension $M=X\times\R$,
this distributional generator admits a one-parameter family of
smooth regularizations
\[
  H_{\mathrm{diss}}^\beta(x,z)
  \;=\; -\gamma\,V(x)\,e^{-\beta z}, \quad \beta>0,
\]
where $V\colon X\to\R_{\ge 0}$ is a smooth Lyapunov function
whose zero set $\{V=0\}$ is the post-fold invariant set $\Gamma$.

\begin{proposition}[Regularization of the fold generator]
\label{prop:regularization}
In the limit $\beta\to\infty$ and $z\to 0^+$ with $\gamma\beta$
fixed, the contact Hamiltonian correction $H_{\mathrm{diss}}^\beta$
reduces to a first-order impulse in the transverse direction,
recovering the momentum jump $p^+-p^-=\mu n$ of the fold map.
\end{proposition}

\begin{proof}[Proof sketch]
At $z=0^+$ and large $\beta$: $e^{-\beta z}\approx 1-\beta z$.
The contact equations give transverse velocity correction
$-\nabla_x H_{\mathrm{diss}}^\beta = \gamma e^{-\beta z}\nabla V
\approx \gamma\nabla V$, which in the normal direction near
$\{V=0\}$ is $\gamma\,\|\nabla V\|\,n = \mu n$ upon identifying
$\mu = \gamma\|\nabla V\|$. As $\beta\to\infty$, the correction
is concentrated at $z=0$, recovering the distributional impulse
of the Heaviside generator $S$.
\end{proof}

\begin{remark}
Proposition~\ref{prop:regularization} should be understood as
a structural correspondence, not an exact limit in the functional-
analytic sense. The precise statement is: the family
$H_{\mathrm{diss}}^\beta$ interpolates between the smooth
contact dynamics (finite $\beta$) and the impulsive fold
map ($\beta\to\infty$). The dm3 framework operates at finite
$\beta$, where all dynamics are smooth.
\end{remark}

The regularized Hamiltonian $H_{\mathrm{diss}}^\beta$ has three
structural properties that make it the correct contact-geometric
continuation of the fold:
\begin{enumerate}[label=(\roman*)]
  \item $H_{\mathrm{diss}}|_\Gamma = 0$: the invariant set
        $\Gamma=\{V=0\}$ is preserved,
  \item $H_{\mathrm{diss}} < 0$ off $\Gamma$: contraction
        toward $\Gamma$ is enforced,
  \item the factor $e^{-\beta z}$ weakens dissipation as
        action $z$ accumulates, encoding the embodiment
        threshold.
\end{enumerate}

Thus the Hamiltonian impulse at $\kappa=\kappa^*$ and the contact
dissipation governing post-fold attraction are two aspects of the
same geometric mechanism: rank-1 injectivity loss followed by
contact stabilization.

%% ------------------------------------------------------------
\subsection{Relation to the dm3 Equations}
\label{subsec:dm3-relation}

Near the stabilized orbit $\Gamma$, Paper~1 (Theorem~C,
\cite{Paper1}) proves that every dm3 system is locally
contact-diffeomorphic to the normal form
\begin{align}
  \dot\rho &= \mu_{\max}(1-e^{-\beta z})\rho + O(\rho^2),
    \label{eq:nf-rho}\\
  \dot\theta &= \omega + O(\rho),
    \label{eq:nf-theta}\\
  \dot z &= \omega - |\mu_{\max}|\rho^2 e^{-\beta z} + O(\rho^3),
    \label{eq:nf-z}
\end{align}
where $\rho$ measures transverse deviation from $\Gamma$,
$\omega=2\pi/T^*$ is the angular frequency, and
$(\mu_{\max},\omega,\beta)$ are the canonical invariant
parameters~\cite{Paper1}.

This normal form can be interpreted directly in terms of the
generative transition sequence $C\to K\to F\to U$:
\begin{itemize}[leftmargin=*]
  \item the approach $\rho\to 0$ under~\eqref{eq:nf-rho}
        corresponds to the unfolding operator $U$: post-fold
        stabilization by gradient descent;
  \item the negativity of $\mu_{\max}$ is the dynamical imprint
        of having passed the geometric threshold $\kappa^*$: the
        fold has occurred and the system is now in the attracting
        regime;
  \item the contact variable $z$ records accumulated dissipation
        generated by folding: the system's history is encoded
        in $z$, and the embodiment threshold $\tau$ measures
        the noise amplitude below which this history stabilizes
        the orbit.
\end{itemize}

In this sense, the dm3 flow is not a separate construction
but the smooth contact-geometric continuation of the fold
dynamics of Volume One. The explicit verification of this
correspondence in the toy model is given in Paper~2~\cite{Paper2},
where the equations~\eqref{eq:nf-rho}--\eqref{eq:nf-z} are
derived by direct Taylor expansion at $r=1$.

%% ------------------------------------------------------------
\subsection{Conceptual Correspondence}
\label{subsec:correspondence-table}

The correspondence established in this section is summarized
in Table~\ref{tab:correspondence}.

\begin{table}[ht]
\centering
\begin{tabular}{ll}
\toprule
\emph{Volume One (geometric)} & \emph{Volume Two / GCM (contact)} \\
\midrule
Curvature threshold $\kappa^*$ & Onset of transverse contraction \\
Fold impulse $p^+-p^-=\mu n$ & Contact Hamiltonian correction
  $H_{\mathrm{diss}}$ \\
Rank-1 Jacobian loss at $F$ & Dissipative contact normal form \\
Unfolding $U$ & Attraction to dm3 orbit $\Gamma$ \\
Symplectic preservation & Contact structure on $M=X\times\R$ \\
$A_1$--$A_3$ singularities & Four dm3 bifurcations (Theorem~C,
  \cite{Paper2}) \\
\bottomrule
\end{tabular}
\caption{Correspondence between the geometric framework of
  Volume One and the contact-geometric framework of
  GCM~\cite{Paper1,Paper2}.}
\label{tab:correspondence}
\end{table}

The fold operator $F$ is thus the geometric progenitor of the
dm3 dynamics: singular in the Hamiltonian picture, regularized
and stabilized in the contact picture.

%% ------------------------------------------------------------
\subsection{Role in the Trilogy}
\label{subsec:trilogy-role}

This section completes the conceptual task left intentionally
open in Volume One: it shows how a purely geometric fold becomes
a dissipative, structurally stable process once contact geometry
is introduced. The next section establishes the quantitative
equivalence between the curvature threshold $\kappa^*$ and the
embodiment threshold $\tau$ — the central bridge theorem of this
volume.
